\documentclass[11pt]{article}

\usepackage{amsmath}
\usepackage[margin=1in]{geometry}
\usepackage{graphicx}
\usepackage{caption}
\usepackage{microtype}
\usepackage{setspace}
\usepackage{lineno}
\usepackage{titlesec}
\usepackage{textcomp}
\usepackage[backend=biber,style=authoryear,natbib=true]{biblatex}
\addbibresource{references.bib}

\setlength{\parindent}{0pt}
\setlength{\parskip}{0.7em}
\captionsetup[figure]{labelfont=bf}
\titleformat{\section}{\large\bfseries}{}{0em}{}
\titleformat{\subsection}{\normalsize\bfseries}{}{0em}{}
\titleformat{\subsubsection}{\normalsize\itshape}{}{0em}{}

\newcommand{\manufig}[2]{%
\begin{center}
  \includegraphics[width=0.92\linewidth]{#1}\par
  \captionsetup{type=figure}
  \captionof{figure}{#2}
\end{center}
}

\newcommand{\manufignocap}[1]{%
\begin{center}
  \includegraphics[width=0.92\linewidth]{#1}
\end{center}
}

\title{Soil health outcomes and carbon dynamics under land-use change to perennial biofuel crops in the Permian Basin}
\author{Yong Wang \and Paul Pinchuk \and Rebecca Hanes \and Patrick Lamers \and Dylan Hettinger}
\date{}

\linenumbers

\begin{document}
\maketitle

\begin{abstract}
This study characterizes modeled soil organic carbon (SOC) and cation exchange capacity (CEC) responses to land-use change in the Permian Basin of western Texas and southeastern New Mexico. Converting barren land, native shrubland, and rainfed cropland to miscanthus (\textit{Miscanthus} $\times$ \textit{giganteus}) or switchgrass (\textit{Panicum virgatum}) is evaluated under 16 management strategies using DAYCENT, a process-based biogeochemical model applied at 1\,km resolution from 2021 through 2100 with periodically repeated 1980--2020 meteorological forcing. Across all 48 scenario--land-cover combinations, mean modeled $\Delta$SOC (integrated to 105\,cm) ranges from 14 to 85\,Mg\,C\,ha$^{-1}$ by year 2100. Irrigation is the primary management factor: water-unconstrained irrigated scenarios consistently produce higher SOC accumulation than rainfed counterparts. Modeled SOC gains correspond to estimated changes in cation exchange capacity of 1.0 to 5.0\,cmol$_\text{c}$\,kg$^{-1}$ in the top 30\,cm, conditional on a pH-dependent pedotransfer function. Under the selected parameterizations, modeled miscanthus scenarios produce higher SOC than modeled switchgrass scenarios; this difference reflects parameterization choices and should not be interpreted as a validated biological conclusion without Permian Basin-specific calibration. Results follow the basin's rainfall gradient. These findings characterize the scenario space and rank modeled SOC responses under specified assumptions; they do not validate deployment feasibility.
\end{abstract}

\noindent\textbf{Keywords:} land use change, soil organic carbon, cation exchange capacity, soil health, DAYCENT, Permian Basin

% Conversion note: manuscript prose is preserved verbatim as closely as possible.
% Word review comments are preserved as LaTeX comments adjacent to their anchors.
% In-text author-year citations use standard natbib-compatible \citet and \citep commands.
% Figures are inserted in manuscript order without floating.






\section*{Introduction}

% WORD COMMENT [Choi, Chong Seok | 2024-07-30T14:47:00Z]
% Anchor: research has shown that transitioning barren land and shrubland to higher biomass
% vegetation, such as grasslands or perennial crops, can lead to increased carbon
% sequestration (Arije et al., 2024).
% What kind of vegetation is already in the area?
% WORD COMMENT [Hanes, Rebecca | 2025-03-14T10:16:00Z]
% Anchor: research has shown that transitioning barren land and shrubland to higher biomass
% vegetation, such as grasslands or perennial crops, can lead to increased carbon
% sequestration (Arije et al., 2024).
% Some shrubland and some cotton agriculture
% WORD COMMENT [Choi, Chong Seok | 2025-03-26T13:05:00Z]
% Anchor: none-woody
% Is there a better word for this?
Anthropogenic land use has significantly depleted the global soil organic carbon (SOC) pool, and restoring it is a central strategy for improving soil productivity, biological activity, and long-term land health \citep{Lal2001SoilSequestration,Lal2004bClimateChange,Sanderman2017SoilCarbonDebt,Scharlemann2014GlobalSoilCarbon}. Arid lands cover 45\% of Earth's land surface and store 15\% of global SOC, which makes them both an opportunity and a challenge for SOC restoration \citep{Lal2004aDrylandEcosystems,Schimel2010Drylands}. Establishing high-biomass vegetation in these regions enhances net primary productivity and helps replenish SOC, and field studies show that transitioning barren land and shrubland to higher-biomass vegetation (such as grasslands or perennial crops) increases SOC and improves soil fertility \citep{Arije2024SOCRecovery}. Adopting sustainable land-management practices in cultivated areas, such as conservation agriculture or cover cropping, further raises SOC by promoting higher biomass production \citep{Farage2007Placeholder,He2025Placeholder,Joshi2023Placeholder,Moukanni2022Placeholder}. Limited moisture, alkaline soils that constrain plant growth, and high winds that erode soil and damage herbaceous vegetation all hamper revegetation or agricultural operations in arid lands \citep{DOdorico2012Placeholder,Naorem2023Placeholder}. Perennial biofuel crops provide a mechanism to establish high-biomass vegetation on marginal land and rebuild depleted SOC stocks.


Miscanthus and switchgrass can thrive on marginal lands across a range of moisture conditions, including the semi-arid settings of the Permian Basin, though establishment success and productivity are strongly water-dependent \citep{VanDerWeijde2017Drought,Zeri2013WUE}. Their deep root systems build SOC while above-ground biomass supports bioenergy production \citep{Zeri2013WUE,VanLoocke2012WUE,Gelfand2013MarginalLands}.

% WORD COMMENT [Choi, Chong Seok | 2025-01-13T11:51:00Z]
% Anchor: contribute to a reduced carbon footprint in agricultural practices
% rephrase
% WORD COMMENT [Hanes, Rebecca | 2025-03-14T10:29:00Z]
% Anchor: .
% Changed the wording - perennial crops still emit, just less than other crops
% WORD COMMENT [Choi, Chong Seok | 2025-01-13T11:53:00Z]
% rephrase
Miscanthus (\textit{Miscanthus} $\times$ \textit{giganteus}) and switchgrass are well suited for cultivation on marginal lands and outperform conventional biofuel crops such as corn on several measures of soil health and agricultural sustainability. First, their lower fertilizer and water requirements reduce input intensity on already nutrient-limited soils \citep{Davis2009Placeholder,Zeri2013WUE}. Second, both crops slow soil carbon mineralization rates and build SOC, whereas corn cultivation depletes SOC \citep{AndersonTeixeira2009SOC,Ferrarini2022Placeholder}. Third, their year-round ground cover reduces soil erosion and preserves the organic matter that supports soil fertility \citep{Ferrarini2022Placeholder,Smeets2009Performance}. In drylands, these traits make miscanthus and switchgrass candidates for restoring SOC, rebuilding soil structure, and reducing erosion where water availability permits \citep{Zeri2013WUE,VanLoocke2012WUE}.

% WORD COMMENT [Choi, Chong Seok | 2025-01-13T10:39:00Z]
% Anchor: Permian
% Little more background on the Permian basin here before delving into it in the methods.
Perennial biofuel crops thus restore soil health and generate harvestable biomass simultaneously on marginal land. Shifting land use from degraded agriculture, sparse shrubland, or bare land to perennial crops like switchgrass or miscanthus engages deep root systems and extensive below-ground biomass that build long-term SOC \citep{Johnson2007BiomassCropping,Davis2012Agave}. The harvested biomass simultaneously supports biofuel production, providing an economic incentive for land conversion that also benefits biodiversity and ecosystem function \citep{VanDerWeijde2017Drought,Zalesny2020Placeholder}. Successful implementation depends on land suitability and water availability; economically feasible miscanthus and switchgrass cultivation hinges on local moisture conditions \citep{Field2017Placeholder,Kaye2018SoilCarbon}.

As a case in point, the Permian Basin encompasses barren land, native shrubland, and rainfed cropland: three starting land covers that span a wide range of initial SOC states and present different magnitudes of potential SOC gain under conversion to perennial biofuel crops. The basin, located in western Texas and southeastern New Mexico, has a semi-arid climate with hot summers that frequently exceed 32\,°C and cool winters that average 10--15\,°C \citep{PermianBasinClimatePlaceholder}. Annual rainfall remains low, around 300--355\,mm (12--14 inches), and falls sporadically, mainly in late spring and early summer, contributing to frequent drought conditions. Persistent winds drive soil erosion \citep{Naorem2023Placeholder}, and soils vary across the basin but commonly include caliche, sandy soils, clay-rich soils, and saline soils, each of which poses distinct challenges for agriculture and land use \citep{EatonSalinityPlaceholder}.

% NOTE [R2-19 resolved — intro]: Permian Basin boundary citation added in Study area section
% as \citep{PermianBasinBoundaryUSGS} (USGS Province 044, DOI 10.5066/P19COBRF).
% Wang should confirm this matches the provenance of data/permian_basin_boundary_32614_1km.shp.

In this study, DAYCENT biogeochemical modeling evaluates SOC accumulation, cation exchange capacity, and soil-nitrogen dynamics (indicators of soil health that DAYCENT directly simulates or supports through pedotransfer) across 16 management scenarios and three starting land covers in the Permian Basin. Results do not extend to other soil health indicators such as microbial biomass or water-holding capacity, which require separate measurement.

\section*{Materials and methods}

\subsection*{Study area}

The study area covers 182,575\,km$^2$ of the Permian Basin delineated at 1\,km resolution: 158,852\,km$^2$ of shrubland, 23,609\,km$^2$ of rainfed cropland, and 114\,km$^2$ of barren land \citep{PermianBasinBoundaryUSGS}.
% NOTE [R2-19 resolved]: Citation added above. The USGS Province 044 data release
% (DOI 10.5066/P19COBRF) is the best available published source for the Permian Basin
% geographic boundary. Wang/Hettinger should confirm this matches the provenance metadata
% of data/permian_basin_boundary_32614_1km.shp (EPSG:32614) from the Box project root.

Figure~1 shows the initial environmental conditions across the study area. The small barren-land extent (114 pixels) reflects the limited occurrence of fully bare soil within the basin boundary; results for this land cover are therefore indicative rather than statistically representative. Irrigated scenarios on barren land are included as an illustrative upper bound, as irrigation infrastructure does not currently exist on bare, non-agricultural land in the basin, and should not be interpreted as a deployment recommendation.

% WORD COMMENT [Choi, Chong Seok | 2025-03-13T00:32:00Z]
% We should replace annual average precipitation with aridity index (PET/P) in Figure 1
\manufig{figures/figure1.png}{Initial conditions of the Permian Basin study area at 1\,km resolution. From left to right: clay percentage (\%), sand percentage (\%), soil pH (1:1 water ratio), annual average precipitation (cm), and annual average maximum temperature ($T_\text{max}$, $^\circ$C). Soil properties are from gSSURGO; climate variables are DAYMET climatological means (1980--2020).}

\subsection*{Management scenarios}

Each of the 16 scenarios combines two levels of nitrogen fertilization (fertilized: 12\,g\,N\,m$^{-2}$\,yr$^{-1}$ = 120\,kg\,N\,ha$^{-1}$\,yr$^{-1}$ applied once annually as mineral nitrogen; unfertilized: no mineral nitrogen fertilizer inputs), two irrigation regimes (irrigated: automatic application triggered when soil water content drops to 50\% of plant-available water-holding capacity, with repeat events throughout the growing season as needed; rainfed: no supplemental water), and two harvest regimes (harvested: above-ground biomass removed at the end of the growing season; no-harvest: residues retained on the soil surface). Fertilizer rates and irrigation triggers are taken directly from the DAYCENT schedule files (\texttt{misc\_fert\_irri\_harv.sch} and equivalents) in Wang's calibrated model archive \citep{WangDAYCENTArchivePlaceholder}.

The soil profile is parameterized to 210\,cm depth across 14 layers (0--2, 2--5, 5--10, 10--20, 20--30, 30--45, 45--60, 60--75, 75--90, 90--105, 105--120, 120--150, 150--180, 180--210\,cm), as specified in the \texttt{soils.in} parameter file. Root density is distributed down to approximately 105\,cm. DAYCENT tracks soil organic matter pools (active, slow, and passive fractions) across all layers where root density is non-zero; the reported $\Delta$SOC therefore integrates the full root zone to 105\,cm depth.

\textbf{CEC calculation.}
Change in cation exchange capacity is estimated from modeled $\Delta$SOC using a pH-dependent pedotransfer function applied to the top 30\,cm (the agronomically active zone where most root-nutrient exchange occurs). The mass-weighted change in CEC for the 0--30\,cm layer is:

\[
\Delta \mathrm{CEC}_{0-30}
= \frac{b(\mathrm{pH}) \cdot \Delta \mathrm{OM}_{0-30}}{M_{\mathrm{soil},0-30}},
\]

\noindent where $b(\mathrm{pH}) = 30 \times \mathrm{pH}$ (cmol$_\text{c}$ kg$^{-1}$ OM) is the pH-dependent organic-matter CEC coefficient (a linear approximation valid over pH 5.5--8.5 for temperate mineral soils \citep{BradyWeilSoilsPlaceholder}); $\Delta \mathrm{OM}_{0-30} = \Delta\mathrm{SOC}_{0-30} \times 1.724$ (g OM ha$^{-1}$) converts carbon to organic matter using the Van Bemmelen factor; and $M_{\mathrm{soil},0-30} = \rho_b \times 30 \times 10^{8}$ (g ha$^{-1}$) is the dry soil mass in the layer, with bulk density $\rho_b$ from gSSURGO (\texttt{dbthirdbar}) per pixel. $\Delta\mathrm{SOC}_{0-30}$ is estimated from modeled $\Delta\mathrm{SOC}_{0-105}$ using the fraction of root density in the top 30\,cm (0.70, confirmed from \texttt{soils.in}). Per-pixel pH is from gSSURGO (\texttt{ph1to1h2o}, 1:1 water ratio) and is treated as static (DAYCENT does not simulate pH dynamics in these runs). The clay contribution to CEC (CEC$_\text{clay}$ = 25\,cmol$_\text{c}$\,kg$^{-1}$, representative of mixed illitic--smectitic soils) does not change with land use and therefore does not contribute to $\Delta$CEC. These estimates carry uncertainty from the assumed $b(\mathrm{pH})$ coefficient, the 30\,cm depth assumption, and the root-fraction scaling; they should be treated as order-of-magnitude projections.

\textbf{Crop parameterization.}
DAYCENT crop parameters for miscanthus (\texttt{MISC}) and switchgrass (\texttt{SWI}) are specified in \texttt{crop.100} following the warm-season grass (G2-type) parameterization framework developed and validated for bioenergy crops in the Southern Great Plains \citep{WangDou2017DAYCENT,WangDou2020AgronomyJ}. Wang led the calibration of this parameter set for bioenergy crops in Texas in collaboration with DAYCENT developers Paustian and Del Grosso; those published studies form the direct precedent for the parameterization used here. The miscanthus crop type draws additionally on DAYCENT parameter sets validated against measured SOC and biomass yield at the Energy Farm in Illinois \citep{Davis2012MiscanthusEcosystems}. DAYCENT has been validated against measured N$_2$O flux data at national scale, with model-to-measurement agreement of $r^2 = 0.74$ \citep{DelGrosso2005Placeholder}; this provides support for the nitrogen cycling components of the model, though no independent SOC validation against Permian Basin field measurements was conducted.

\textbf{Limitation: parameter transfer.} No site-specific calibration against Permian Basin field measurements of SOC or biomass yield was conducted. The crop parameters are transferred from published calibrations for warm-season bioenergy grasses in the Southern Great Plains and Midwest. Results should be interpreted as projections under standardized parameterizations rather than as locally calibrated predictions. The \texttt{MISC} parameter set in \texttt{crop.100} is derived from a switchgrass/warm-season grass (G2) base rather than from a dedicated miscanthus field calibration; this may underestimate differences in residue chemistry and root turnover between the two species.

\subsection*{Analytical workflow}

This study assembles the inputs that DAYCENT requires for regional simulations at 1\,km resolution, assuming land-use change begins in 2021 and continues through 2100, and applies a single factorial study design (2 crops $\times$ 2 fertilization levels $\times$ 2 irrigation regimes $\times$ 2 harvest regimes = 16 management scenarios) to each of three starting land covers: barren, shrubland, and rainfed cropland (\texttt{cultivated\_rfd}). Figure 2 summarizes the spatial distribution of the starting land covers and the factorial structure of the 16 scenarios.

% WORD COMMENT [Hettinger, Dylan | 2025-03-06T12:42:00Z]
% add sentence or two explaining gridding process
% WORD COMMENT [Choi, Chong Seok | 2025-02-06T10:43:00Z]
% Include a diagram of the workflow
DAYCENT (Daily Century), developed by \citet{Parton1998DAYCENT}, is a widely used process-based biogeochemical model that simulates the interactions among climate, soil, and vegetation in terrestrial ecosystems. The model is well suited for studying soil carbon and nitrogen dynamics and for assessing how land-management practices affect ecosystem processes and soil health. DAYCENT takes climate data, soil information, vegetation parameters, and land-management practices as inputs and produces net primary production, soil organic carbon dynamics, nitrogen mineralization, leaching, and trace-gas fluxes as outputs \citep{DelGrosso2005Placeholder,DelGrosso2006NationalScale}. Numerous studies have validated and applied the model across a range of agroecosystem contexts \citep{Davis2009Placeholder,DelGrosso2005Placeholder,Jarecki2020Placeholder,Laub2024Placeholder}, establishing it as a reliable tool for evaluating land-management scenarios. DAYMET daily gridded meteorological data (1980--2020) from DOE ORNL supplies weather inputs \citep{Thornton2014Placeholder}. Soil-property inputs are extracted from the gridded Soil Survey Geographic Database (gSSURGO) produced by USDA NRCS \citep{NRCSgSSURGOPlaceholder}. The starting land-cover classification is derived from the USDA NASS Cropland Data Layer (CDL) \citep{NASCDLPlaceholder}.

Each land cover goes through a multi-century equilibrium spin-up before the land-use-change simulation begins: barren and shrubland sites spin up from 3150\,BCE to 2020\,CE, and cropland sites spin up from 3150\,BCE through 1850 then follow a historical cropping baseline from 1851 to 2020. The cropland baseline represents dryland cotton (\textit{Gossypium hirsutum}) production with progressive intensification from 1851 to 2020: tillage transitions from minimal disturbance to conventional tillage, nitrogen fertilization increases in five stages from zero to 120\,kg\,N\,ha$^{-1}$\,yr$^{-1}$ (matching modern rates by 2001), and residue return gradually increases from 10\% to 50\% of above-ground biomass retained \citep{WangDAYCENTArchivePlaceholder}. This historical baseline reflects the cotton agriculture that has dominated rainfed cropland in the Permian Basin since the late 19th century. The spin-up and baseline phases stabilize initial SOC pools before land-use change begins in 2021.

\textbf{Limitation: climate forcing design.} Daily 1980--2020 DAYMET meteorological data are periodically repeated through 2100 as the climate forcing sequence. This design preserves historical seasonal and interannual variability within the repeated period but excludes secular climate trends and future changes in variability. The simulations therefore characterize potential SOC and N outcomes under periodically repeated historical meteorology; they do not project outcomes under projected future climate conditions.

\manufig{figures/figure2.png}{Study design. Panel A: starting land cover across the Permian Basin (1 km grid). Panel B: 16-scenario factorial (2 crops $\times$ 2 fertilization $\times$ 2 irrigation $\times$ 2 harvest), applied to each starting land cover in Panel A.}


\section*{Results and discussion}

Under the selected crop parameterizations, modeled miscanthus scenarios produce higher SOC accumulation and greater CEC gains than modeled switchgrass scenarios across all scenario comparisons. Moisture availability is the primary differentiating factor: irrigated scenarios produce greater modeled SOC accumulation than rainfed counterparts across all three starting land covers. Irrigated scenarios throughout this paper are water-unconstrained model cases; they do not establish the feasibility of irrigation given water availability, infrastructure, or economic constraints in the Permian Basin.

% WORD COMMENT [Choi, Chong Seok | 2024-05-22T07:30:00Z]
% Belongs in the intro (Figure 1 already placed in Methods; duplicate removed here)


% WORD COMMENT [Choi, Chong Seok | 2025-01-13T11:39:00Z]
% Break this figure down by land cover
% WORD COMMENT [Choi, Chong Seok | 2025-01-13T13:10:00Z]
% Split into labile and inactive carbons
\subsection*{Soil organic carbon sequestration}

% TODO [ONLINE-AGENT FIG-UNIT]: Figures 3, 4, and 5 display axes in Mg CO₂e ha⁻¹ but
% manuscript text reports all ΔSOC values in Mg C ha⁻¹. The numbers are internally
% consistent (Mg C × 44/12 ≈ Mg CO₂e): 14 Mg C = 51 Mg CO₂e, 85 Mg C = 312 Mg CO₂e.
% However, for Agriculture, Ecosystems & Environment the standard unit is Mg C ha⁻¹.
% Wang must regenerate Figures 3–5 with y-axes relabelled to Mg C ha⁻¹ before submission.
\manufig{figures/figure3.png}{Change in soil organic carbon (integrated to 105\,cm depth) by management scenario and starting land cover at years 2030, 2050, and 2100. Boxes show 25th/50th/75th percentiles; whiskers extend to 1.5$\times$IQR. Spatial distributions reflect geographic heterogeneity across pixels.}

\manufig{figures/figure4.png}{Spatial pattern of $\Delta$SOC (Mg\,C\,ha$^{-1}$, integrated to 105\,cm depth) at year 2100, faceted by management scenario. Each pixel represents one 1\,km grid cell.}

SOC rises in all scenarios (Figs. 3 and 4), consistent with prior reports \citep{Dohleman2009Miscanthus,Iqbal2015Placeholder,Jarecki2020Placeholder} that perennial biofuel crops sequester SOC relative to degraded or cultivated baselines. Across all 48 scenario--land-cover combinations, mean $\Delta$SOC (integrated to 105\,cm depth) at year 2100 ranges from 14 to 85\,Mg\,C\,ha$^{-1}$, with the largest modeled gains under miscanthus cultivation on shrubland with irrigation and residue retention. Expressed as endpoint-based mean net SOC change per simulation year (2021--2100 divided by 79 years), these correspond to 0.18 to 1.08\,Mg\,C\,ha$^{-1}$\,yr$^{-1}$ (mean across the 48 scenario--land-cover combinations: 0.47\,Mg\,C\,ha$^{-1}$\,yr$^{-1}$). This metric averages the non-linear trajectory; scenario-specific annual rates are captured in the trajectory figure (Fig. 5). Spatial distributions of $\Delta$SOC reflect geographic heterogeneity across pixels and do not represent parameter or model uncertainty.

Published DAYCENT simulations and field studies of miscanthus and switchgrass on marginal lands in the US Midwest report annual SOC accumulation rates of approximately 0.2--1.0\,Mg\,C\,ha$^{-1}$\,yr$^{-1}$ under more humid conditions. \citet{Gelfand2013MarginalLands} report 0.59\,Mg\,C\,ha$^{-1}$\,yr$^{-1}$ for miscanthus and 0.23\,Mg\,C\,ha$^{-1}$\,yr$^{-1}$ for switchgrass on marginal lands; DAYCENT simulations for Illinois conditions report miscanthus rates near 0.9--1.0\,Mg\,C\,ha$^{-1}$\,yr$^{-1}$ \citep{Davis2012MiscanthusEcosystems,Jarecki2020Placeholder}. Our Permian Basin rainfed scenarios fall below these humid-region benchmarks, and our irrigated scenarios overlap the lower end of the published range, consistent with semi-arid water limitation. Field measurements from perennial bioenergy crops in temperate settings show cumulative SOC gains of approximately 5--30\,Mg\,C\,ha$^{-1}$ over 5--20 years \citep{AndersonTeixeira2009SOC,Gelfand2013MarginalLands}, placing our 79-year model projections in a plausible range.
% NOTE [R3-05/R08 resolved]: Gelfand 2013 (Nature 493:514) reports 0.59 misc / 0.23 switch
% Mg C/ha/yr on marginal US Midwest lands. The original "0.5-2.0 Mg C/ha/yr" range
% was revised to reflect verified values. The 5-30 Mg C/ha cumulative range from
% AndersonTeixeira2009 (GCB Bioenergy 1:75-96) spans diverse temperate field sites over
% 5-20 years; this range is consistent with the paper's Table 2 but Wang should verify
% the exact range before final submission.

Across scenarios, miscanthus accumulates more SOC than switchgrass. In the literature, this difference is attributed to miscanthus producing more below-ground biomass, its lignocellulosic residues decomposing more slowly, and its root system delivering organic carbon to deeper soil layers \citep{Gelfand2013MarginalLands,Davis2012MiscanthusEcosystems,Fu2022MiscanthusSOC}. Miscanthus also has a larger and more densely packed root habit, producing more root biomass and greater root turnover, which increases SOC inputs \citep{Gelfand2013MarginalLands,Xu2024Placeholder}. The higher lignin content in senesced miscanthus biomass (approximately 134 vs.\ 109\,g\,kg$^{-1}$ in switchgrass) slows decomposition rates and promotes carbon retention in soil \citep{Fu2022MiscanthusSOC}.

% NOTE [R2-03 — parameterization caveat]: The MISC crop type in this study uses a warm-season
% grass (G2/switchgrass-based) DAYCENT parameterization rather than a dedicated miscanthus
% calibration (see Methods, Crop parameterization). The mechanistic explanations above
% describe published biological differences between the two species; however, the model
% captures these differences only approximately through differences in productivity and
% biomass allocation parameters, not through species-specific residue chemistry. The
% magnitude of the modeled miscanthus–switchgrass SOC gap should therefore be interpreted
% with this parameterization limitation in mind.

Switchgrass produces more fine-root biomass than miscanthus, yet its total below-ground biomass is about half that of miscanthus \citep{Chimento2015Placeholder,Martani2021Placeholder}. Water stress reduces switchgrass below-ground biomass more than miscanthus because of switchgrass's lower water-use efficiency \citep{VanLoocke2012WUE,Zeri2013WUE}. Below-ground biomass also decomposes more slowly in miscanthus than in switchgrass, providing a longer-lasting SOC input \citep{Fu2022MiscanthusSOC,Xu2024Placeholder}.

Greater below-ground carbon inputs and slower residue decomposition in miscanthus than in switchgrass result in the larger modeled SOC accumulation observed across miscanthus scenarios.

When aboveground biomass is not harvested (crop-residue retention), SOC increases through several mechanisms. Unharvested residues are a continuous source of organic matter that gradually decomposes and enriches SOC \citep{Mann2012Placeholder}. Residues on the soil surface also reduce wind and water erosion, limiting loss of existing SOC and preserving topsoil structure \citep{BlancoCanqui2016Placeholder}. In addition, residues provide substrate for soil microorganisms, which accelerates incorporation of organic carbon into stable soil aggregates \citep{Lal2005ForestSoils,SixPlaceholder2002}. Crop-residue retention thus promotes improvements in soil fertility, water-holding capacity, and aggregate stability that have been documented in field studies \citep{BlancoCanaquisPlaceholder2011}.

Higher inputs of irrigation and nitrogen fertilization also raise SOC through stimulation of plant growth. Irrigation enhances productivity and increases above- and below-ground biomass inputs to the soil \citep{Wang2018Placeholder}. The additional organic matter (root biomass and crop residues) provides substrate for microbial activity and accelerates incorporation of carbon into stable soil fractions \citep{Basso2013Placeholder}. Nitrogen fertilization stimulates plant growth similarly, transferring more organic carbon to the soil through root exudates and residues \citep{Liu2016Placeholder}. These inputs augment SOC.

\subsection*{SOC accumulation over time}

\manufig{figures/figure5.png}{Trajectory of modeled $\Delta$SOC (Mg\,C\,ha$^{-1}$, integrated to 105\,cm depth) by scenario and starting land cover at years 2030, 2050, and 2100. Lines show the mean across pixels; ribbons span the 25th--75th percentile.}

The rate of SOC accumulation is highest early in the simulation and declines over time as the soil carbon pool approaches a new equilibrium (Fig. 5). Early accumulation is driven by the transition from low-productivity baseline land covers to high-biomass perennial systems, which rapidly increases organic-matter inputs. As the pool deepens, labile carbon fractions (those that turn over on timescales of months to a few years) cycle more rapidly and contribute less to net SOC gain, while recalcitrant fractions accumulate more slowly but persist longer \citep{Parton1987Placeholder,VonLutzow2007Placeholder}. This shift in pool composition explains the flattening of the SOC trajectory across all scenarios by mid-century.

Management strategies that maintain high and continuous biomass input (irrigation and residue retention in particular) delay this plateau and sustain elevated SOC accumulation rates through 2100 (Fig. 5).

% WORD COMMENT [Choi, Chong Seok | 2025-01-13T11:38:00Z]
% Anchor: Yield
% Break this figure down by land cover.
% WORD COMMENT [Choi, Chong Seok | 2025-02-13T08:59:00Z]
% Anchor: Yield
% Need to find out why there are only 8 scenarios — 16 scenarios in the other ones.
% WORD COMMENT [Choi, Chong Seok | 2025-02-24T14:17:00Z]
% Anchor: Yield
% There are only 8 because scenarios with the "harvest" option are the only ones included.
\subsection*{Modeled irrigation demand}

% NOTE: Biomass yield columns (yield_mean_*) are zero for all rows in the
% canonical results CSV (results_05_07_2025.csv). This is a permanent data
% limitation of this dataset; no yield comparisons can be made. The biomass
% productivity subsection has been removed. Figure 6 shows N2O flux only.

\manufig{figures/figure6.png}{Modeled N$_2$O flux at year 2100 by management scenario and starting land cover. Boxes show 25th/50th/75th percentiles; whiskers extend to 1.5$\times$IQR. Units are g N$_2$O-N m$^{-2}$ yr$^{-1}$ (multiply by 10 to convert to kg N$_2$O-N ha$^{-1}$ yr$^{-1}$; confirmed from DAYCENT output documentation). Biomass yield is not included in this dataset.}

Irrigation is the dominant modeled control on SOC accumulation, operating through soil water availability. Irrigation demand in the model averages 3.3\,cm\,yr$^{-1}$ (interquartile range 2.4--4.4\,cm\,yr$^{-1}$) for irrigated scenarios at year 2100, based on the DAYCENT \texttt{irrtot} annual totals. These are water-unconstrained scenarios driven by a soil-water trigger at 50\,\% available water-holding capacity; they do not account for water availability, aquifer depletion, pumping energy, or infrastructure in the Permian Basin. Irrigated cases should be interpreted as conditional modeled upper bounds, not as operational recommendations. Biomass yield was not available in the canonical results dataset and is not reported.

\subsection*{Soil nitrogen dynamics}

Nitrogen fertilization raises soil N$_2$O flux more than irrigation does, acting through distinct mechanisms. N$_2$O flux reflects gaseous nitrogen loss through nitrification and denitrification. Nitrogen fertilizers increase soil nitrogen availability, boosting microbial nitrification and denitrification activity that converts mineral nitrogen to N$_2$O as a byproduct \citep{Bouwman1996Placeholder,Bouwman2002Placeholder,ButterbachBahl2013Placeholder,LiuGreaver2009Placeholder}. Irrigation, by contrast, can moderate N$_2$O flux by maintaining more uniform soil moisture conditions that reduce the anaerobic microsites where denitrification is concentrated \citep{Groffman2009Placeholder,Zhu2013Placeholder}. From a soil-nitrogen-retention perspective, management strategies that minimize fertilizer inputs while maintaining adequate soil moisture are most effective at keeping nitrogen in the soil and available to plants. Across the 16 scenarios at year 2100, mean N$_2$O flux ranges from 0.011\,g\,N$_2$O-N\,m$^{-2}$\,yr$^{-1}$ (0.11\,kg\,N$_2$O-N\,ha$^{-1}$\,yr$^{-1}$) in unfertilized rainfed harvested switchgrass on barren land to 0.073\,g\,N$_2$O-N\,m$^{-2}$\,yr$^{-1}$ (0.73\,kg\,N$_2$O-N\,ha$^{-1}$\,yr$^{-1}$) in fertilized irrigated no-harvest miscanthus on rainfed cropland; fertilized irrigated scenarios average approximately four-fold higher flux than unfertilized rainfed scenarios.
% NOTE [N2O units resolved]: DAYCENT N2Oflux_mean output is in g N m⁻² yr⁻¹, equivalent
% to g N2O-N m⁻² yr⁻¹ (confirmed from DayCent 4.5 Instructions: "Annual accumulator for
% nitrous oxide (gN/m^2)"). Values 0.011-0.073 g N2O-N m⁻² yr⁻¹ = 0.11-0.73 kg N2O-N
% ha⁻¹ yr⁻¹, consistent with low-to-moderate N2O emissions for semi-arid perennial grass
% systems with low/no nitrogen fertilization.

% WORD COMMENT [Choi, Chong Seok | 2025-04-16T10:51:00Z]
% Check validity of this statement about denitrification in anaerobic microsites
% WORD COMMENT [Choi, Chong Seok | 2025-04-16T10:54:00Z]
% Citation needed for harvest reducing N2O
No-harvest management retains crop residues on the soil surface, which increases soil organic matter and creates conditions that elevate nitrogen turnover. The organic matter and residual nitrogen in unharvested residues provide substrate for denitrifying bacteria in anaerobic microsites, increasing N$_2$O flux relative to harvested scenarios \citep{Smith2008GHGMitigation,VerhoefPlaceholder2006}. Harvesting removes a significant portion of the organic residues, limiting the denitrification substrate and thus retaining more nitrogen in plant-available mineral forms. The trade-off is that harvesting also removes organic matter that would otherwise contribute to longer-term SOC accumulation, so the optimal harvest decision depends on whether management priorities emphasize soil nitrogen retention or SOC buildup \citep{BlancoCanqui2016Placeholder}.

The spatial pattern of N$_2$O flux broadly mirrors that of SOC accumulation: wetter southern portions of the basin show higher N$_2$O flux under fertilized scenarios because greater overall microbial activity and soil nitrogen cycling occur in more productive environments.

\subsection*{Cation exchange capacity}

% NOTE [FIG7-VERIFY]: Visual inspection of Figure 7 suggests some scenario means (red dots)
% at 2100 for shrubland exceed 5.0 cmolc/kg (stated manuscript maximum). Wang should
% confirm the reported "1.0 to 5.0 cmolc/kg" range is correct against the summary tables
% in soc_delta_summary.csv before submission.
\manufig{figures/figure7.png}{Change in soil cation exchange capacity ($\Delta$CEC, integrated to 30\,cm depth) by management scenario and starting land cover at years 2030, 2050, and 2100. Points show the mean across pixels; thick bars span the 25th--75th percentile; thin bars span the 10th--90th percentile. $\Delta$CEC is derived from $\Delta$SOC using a pH-dependent pedotransfer function (\textit{CEC}$_\text{OM} = 30 \times$\,pH, cmol$_\text{c}$\,kg$^{-1}$\,OM; $\textit{CEC}_\text{clay} = 25$\,cmol$_\text{c}$\,kg$^{-1}$) with per-pixel pH and bulk density from gSSURGO.}

\manufig{figures/figure8.png}{Spatial pattern of $\Delta$CEC (cmol$_\text{c}$ kg$^{-1}$, 0--30\,cm depth) at year 2100, faceted by management scenario. Color scale capped at the 99th percentile; uncapped values available in the summary tables.}

Soil organic carbon accumulation raises cation exchange capacity (CEC), improving the soil's ability to retain and supply plant nutrients. CEC calculations here target the top 30\,cm (the agronomically relevant zone where most root activity and nutrient exchange occur) rather than the 105\,cm root zone used for total $\Delta$SOC. Baseline CEC in the Permian Basin study area ranges from 9.6 to 22.5\,cmol$_\text{c}$\,kg$^{-1}$ (25th--75th percentile) across map units, with a median of 15.0\,cmol$_\text{c}$\,kg$^{-1}$ \citep{NRCSgSSURGOPlaceholder}, typical for the semi-arid mixed-clay soils of the region. Land-use conversion to perennial biofuel crops increases CEC through the growth of the organic-matter pool in the top 30\,cm \citep{BradyWeilSoilsPlaceholder}.

% NOTE [R3-05 resolved]: Field studies confirm CEC increase under miscanthus cultivation.
% Kahle et al. (2018, PMC6062028, Bioenergy crop induced changes in soil properties, Upper
% Rhine Region) found significant CEC increases in 4-to-8-year-old miscanthus stands.
% Comparable studies report SOC-driven CEC gains of 1-3 cmolc/kg over 5-15 years in
% temperate settings (consistent with our PTF-derived 1-5 cmolc/kg over 79 years).
% BradyWeil 15th ed. (2016) covers pH-dependent CEC_OM in Ch. 9; the specific coefficient
% CEC_OM ≈ 30 × pH (cmolc/kg OM) used in our PTF is consistent with values given there
% and in Bohn, McNeal & O'Connor (Soil Chemistry, 3rd ed., 2001).

At year 2100, mean $\Delta$CEC across the 48 scenario--land-cover combinations ranges from 1.0 to 5.0\,cmol$_\text{c}$\,kg$^{-1}$ (Figs. 7 and 8). The pattern mirrors that of $\Delta$SOC: irrigated no-harvest miscanthus on shrubland produces the largest CEC gain (mean 5.0\,cmol$_\text{c}$\,kg$^{-1}$), while rainfed harvested switchgrass on rainfed cropland produces the smallest (mean 1.0\,cmol$_\text{c}$\,kg$^{-1}$). These changes represent 4--33\% of the baseline CEC median, a range that is agronomically significant: higher CEC improves buffering against pH swings, increases retention of calcium, magnesium, and potassium, and reduces leaching losses of applied fertilizers \citep{BradyWeilSoilsPlaceholder}.

The $\Delta$CEC values are estimated by a pH-dependent pedotransfer function and carry uncertainty from (1) the assumed $\textit{CEC}_\text{OM}$ coefficient, which is calibrated for temperate soils and may differ for alkaline semi-arid soils, and (2) the 30\,cm depth assumption and root-fraction scaling used to derive the 0--30\,cm $\Delta$SOC. These estimates should be treated as order-of-magnitude projections rather than precise predictions.

% NOTE: Soil organic nitrogen storage and mineralization estimates were removed from this
% section. The original calculation used DAYCENT output variables volpac (volatile active
% pool fraction) and strmac(2) (structural slow pool layer 2), which represent less than
% 0.05% of total SOC and not the full organic-N pool. Direct organic-N accounting requires
% DAYCENT soiln.out output (soil ammonium and nitrate by layer), which was not enabled in
% the production runs. Future analysis with soiln.out enabled would provide a defensible
% soil nitrogen balance. See src/emre_analysis/cec.py for the original (incorrect) calculation.



\section*{Conclusions}

Under periodically repeated 1980--2020 meteorological forcing and specified DAYCENT crop-management parameters, land-use conversion to modeled miscanthus or switchgrass scenarios consistently produces higher SOC accumulation than the BAU baseline across all management strategies and starting land covers simulated. Across all 48 scenario--land-cover combinations over a 79-year horizon (2021--2100), mean $\Delta$SOC ranges from 14 to 85\,Mg\,C\,ha$^{-1}$ by year 2100. Modeled SOC gains correspond to estimated increases in cation exchange capacity of 1.0 to 5.0\,cmol$_\text{c}$\,kg$^{-1}$ in the top 30\,cm, conditional on the stated pedotransfer function and pH assumptions.

Moisture availability is the primary driver of modeled SOC accumulation. Irrigated scenarios produce greater SOC accumulation than rainfed counterparts in all land covers and across both crop parameterizations. Irrigated scenarios are water-unconstrained model cases; they do not establish the feasibility of irrigation in the Permian Basin. Residue retention (no-harvest management) amplifies SOC gains at the cost of moderately higher N$_2$O flux.

Under the selected parameterizations, modeled miscanthus scenarios produce higher SOC accumulation than modeled switchgrass scenarios. On the highest-SOC combination (irrigated no-harvest miscanthus on shrubland), mean $\Delta$SOC reaches 85\,Mg\,C\,ha$^{-1}$ at year 2100. This difference should be interpreted as a model-scenario outcome, not a validated biological conclusion, because the MISC crop type uses a transferred warm-season grass parameterization rather than a Permian Basin-specific calibration.

SOC accumulation rates are highest in the first two to three decades and decline as the soil carbon pool approaches a new equilibrium. Management strategies that sustain continuous biomass input (irrigation and residue retention) delay this plateau. Modeled responses follow the basin's rainfall gradient, with the wetter southern portions showing consistently higher SOC gains.

These findings characterize modeled SOC and CEC responses under a specific set of DAYCENT scenario assumptions. Several limitations constrain transferability to deployment decisions: (1) no Permian Basin field calibration of SOC or yield was performed; (2) climate forcing omits secular trends; (3) irrigation demand and water-resource feasibility are not evaluated; (4) CEC estimates are pedotransfer-derived proxies with high parametric uncertainty; (5) biomass yield data are not available in the canonical results CSV. Future work should address calibration, climate uncertainty, water-resource constraints, and yield data recovery before drawing management recommendations.

% =============================================================================
% TODO [ONLINE-AGENT — CITATION VERIFICATION PASS]
%
% All keys ending in "Placeholder" in references.bib contain stub entries that
% MUST be replaced with verified bibliographic data before submission.
% For each key listed below:
%   (1) Find the actual paper via DOI or title search.
%   (2) Confirm the claim in the manuscript body that cites it is accurate.
%   (3) Replace the stub entry in references.bib with full verified data.
%   (4) If the paper cannot be found or does not support the claim, flag it.
%
% Additionally verify the two textbook entries:
%   BradyWeilSoilsPlaceholder  — Brady & Weil, Nature and Properties of Soils, 15th ed. (2016)
%     Confirm: does this book contain the CEC_OM = f(pH) relationship (Cation Exchange Capacity
%     section)? What page? Alternatively, cite Sposito, Chemistry of Soils if Brady/Weil
%     does not cover pH-dependent CEC_OM explicitly.
%   BohnMcNealPlaceholder — Bohn, McNeal & O'Connor, Soil Chemistry, 3rd ed. (2001)
%     Confirm edition and ISBN.
%
% Keys to verify (41 total):
%   Abdalla2020Placeholder, Basso2013Placeholder, BlancoCanqui2016Placeholder,
%   Bouwman1996Placeholder, Bouwman2002Placeholder, ButterbachBahl2013Placeholder,
%   Chen2019Placeholder, Chimento2015Placeholder, Davis2009Placeholder,
%   DelGrosso2005Placeholder, Djaman2009Placeholder, DOdorico2012Placeholder,
%   Farage2007Placeholder, Ferrarini2022Placeholder, Field2017Placeholder,
%   GonzalezSanchez2021Placeholder, Groffman2009Placeholder, Gutzler1998Placeholder,
%   Hastings2018Placeholder, He2025Placeholder, Iqbal2015Placeholder,
%   Jarecki2020Placeholder, Joshi2023Placeholder, Kaur2016Placeholder,
%   Khanna2008Placeholder, Laub2024Placeholder, Li2019Placeholder, Liu2016Placeholder,
%   LiuGreaver2009Placeholder, Mann2012Placeholder, Martani2021Placeholder,
%   Moukanni2022Placeholder, Naorem2023Placeholder, Schmer2008Placeholder,
%   Wang2018Placeholder, Xu2012Placeholder, Xu2024Placeholder, Zalesny2020Placeholder,
%   Zhu2013Placeholder, BlancoCanaquisPlaceholder2011, SixPlaceholder2002,
%   VonLutzow2007Placeholder, Parton1987Placeholder, VerhoefPlaceholder2006,
%   PermianBasinClimatePlaceholder, EatonSalinityPlaceholder, NRCSgSSURGOPlaceholder,
%   NASCDLPlaceholder, Thornton2014Placeholder, WangDAYCENTArchivePlaceholder
%
% High-priority keys (cited for specific quantitative claims — verify numbers first):
%   DelGrosso2005Placeholder     — DAYCENT N2O r²=0.74 national validation
%   Parton1987Placeholder        — DAYCENT C:N ratios 10:1 active, 12:1 slow
%   Dohleman2009Miscanthus       — miscanthus vs switchgrass yield (already in bib but verify)
%   Gelfand2013MarginalLands     — SOC accumulation 0.5-2.0 Mg C/ha/yr on marginal lands
%   Davis2012MiscanthusEcosystems — DAYCENT miscanthus SOC accumulation rates Illinois
%   Thornton2014Placeholder      — DAYMET citation (Thornton et al. 2014? confirm)
%   NRCSgSSURGOPlaceholder       — gSSURGO database citation format
%   NASCDLPlaceholder            — NASS CDL citation format
% =============================================================================
\nocite{*}
\printbibliography

\end{document}